\begin{table*}[t]
\centering
\caption{代表性方向的能力比较。符号“部分”表示论文可能支持该能力，但未同时给出完整网络、板端逐节点整数一致性与跨 pass 调度证据。}
\label{tab:related}
\small
\resizebox{0.98\textwidth}{!}{%
\begin{tabular}{lccccc}
\toprule
方向/代表工作 & 层间映射 & 外部HWC协同 & 跨pass重放 & 完整检测DAG & 逐节点板端证据\\
\midrule
固定数据流 ASIC（Eyeriss/TPU） & 部分 & 否/内部格式 & 否 & 否 & 否 \\
可重构数据流（MAERI/Simba） & 是 & 部分 & 部分 & 否 & 否 \\
FPGA 生成框架（FINN/fpgaConvNet） & 是 & 部分 & 部分 & 部分 & 部分 \\
YOLO FPGA 加速器 & 部分 & 部分 & 部分 & 是/不一 & 少见 \\
FEATHER & 是 & 是 & 非本文语义 & 非重点 & 否 \\
本文 LASA & 是 & 是 & 是 & 是 & 是 \\
\bottomrule
\end{tabular}
}
\end{table*}
